\subsection{The bounded forms-through-points reconstruction}
\label{subsec:ftb-bounded-flat-projective-reconstruction}

We use the source's stronger density condition.
A set \(P\) of prime ideals of \(R\) is \emph{very dense} if
\[
 \forall q\in\operatorname{Spec}R,\ \forall f\notin q,\
 \exists p\in P:\quad q\subseteq p,\quad f\notin p.
 \tag{M6}\label{eq:ftb-very-dense-primes}
\]
Thus every nonempty relative basic open in every irreducible
closed subset meets \(P\).
Ordinary density in \(\operatorname{Spec}R\) is not enough
for the argument below.

For a finite set \(\Sigma\), write \(A_e(R)=R[X_i:i\in\Sigma]_e\)
for the finite free module of homogeneous forms of degree \(e\).
For finitely many vectors \(p_h\in R^\Sigma\), indexed by a
finite set \(\mathcal H\), set
\[
 E_d:A_d(R)\longrightarrow R^{\mathcal H},\quad
             F\longmapsto(F(p_h))_h,\qquad
             J_d(p)=(\ker E_d)\,R[X_i].
 \tag{M7}\label{eq:ftb-bounded-vanishing-ideal}
\]
In particular \(J_d\) is generated in the single degree \(d\);
it is not the ideal of all forms vanishing at the points.

For a field map \(x:R\to K\), a homogeneous ideal \(I_x\)
has \(\mathrm{CutBy}(Q,d,p_x)\) if, for every \(e\ge d\),
\[
 \dim_K(I_x)_e=Q(e),\qquad
       (I_x)_e=(J_d(xp))_e.
 \tag{M8}\label{eq:ftb-bounded-cutby}
\]
The polynomial \(Q\in\mathbb Q[T]\) therefore records the
dimension of the \emph{ideal} part, not the quotient part.
The integer \(d\), the point labels and \(Q\) are the same
in every fibre used below.

\begin{srlemma}[Very dense constant rank over a reduced ring]
\label{lem:ftb-very-dense-projective}
Let \(R\) be reduced, \(P\) very dense, and \(M\) a finite
\(R\)-module.
If \(\dim_{k(p)}(M\otimes_Rk(p))\) has the same value for
all \(p\in P\), then \(M\) is finite projective.
\end{srlemma}
\begin{proof}
At an arbitrary prime \(q\), finite generation gives a
basic open neighborhood on which fibre dimension is at
most its value at \(q\).
Choose \(p\in P\) above \(q\) in this neighborhood by
\eqref{eq:ftb-very-dense-primes}.
Dimension cannot decrease on specialization from \(q\)
to \(p\).
The two inequalities force the fibre dimension at \(q\)
to equal the prescribed constant.
Thus the fibre dimension is locally constant everywhere.

We spell out the reduced-ring step, which also proves
finite presentation when \(R\) is not Noetherian.
Lift a basis of the fibre at \(q\) to a map
\(\varphi:R^r\to M\).
After shrinking by \(f\notin q\), this map is fibrewise
surjective and every fibre has dimension \(r\).
For a relation \(v\in\ker\varphi\) and a prime \(l\) not
containing \(f\), the fibre map is an isomorphism, so
every coordinate of \(v\) belongs to \(l\).
If \(l\) contains \(f\), it contains every coordinate of
\(fv\) as well.
Thus each coordinate of \(fv\) belongs to every prime
of \(R\); reducedness gives \(fv=0\).
After inverting \(f\) the map is injective.
Its cokernel is finite with zero fibre at every prime
in that open, so it is zero there.
Hence \(M_f\simeq R_f^r\).
These basic opens cover the spectrum.
A finite subcover supplies finite presentation, and a
finitely presented locally free module is projective.
\end{proof}

\begin{srproposition}[The actual flat projective family from bounded torsion forms]
\label{prop:ftb-flat-projective-from-bounded-forms}
Let \(R\) be reduced, \(\Sigma,\mathcal H\) finite, and
let \(p_h\in R^\Sigma\), \(d\in\mathbb N\),
\(Q\in\mathbb Q[T]\) be fixed.
Suppose a very dense set \(P\) has, for every \(p\in P\),
a field \(K\), a map \(x:R\to K\) with kernel \(p\), and
a homogeneous ideal \(I_x\) satisfying
\eqref{eq:ftb-bounded-cutby}.
Then \(J=J_d(p)\) is finitely generated, and
\[
 \begin{split}
 &(A_e(R)/J_e)\text{ is finite projective for every }e\ge d,\\
 &J\,K[X_i]=J_d(xp)\quad\text{for every field map }x:R\to K.
 \end{split}
 \tag{M9}\label{eq:ftb-bounded-form-projectivity}
\]
There is a projective, flat, finitely presented scheme
\(Z/R\), with its actual coordinate line and sections,
representing the homogeneous ideal \(J\).
For every further field map \(x:R\to K\) and every homogeneous
ideal \(I_x\) satisfying \eqref{eq:ftb-bounded-cutby}, the
coordinate representation of \(I_x\) is isomorphic to the
actual fibre of \(Z\), compatibly with the line and all
section-coordinate vectors.
The same \(d,Q,p_h\) occur throughout.
\end{srproposition}
\begin{proof}
Consider \(E_d\) in \eqref{eq:ftb-bounded-vanishing-ideal}.
At a good prime, the degree-\(d\) clause in
\eqref{eq:ftb-bounded-cutby} identifies
\(\ker(E_d\otimes K)\) with \((I_x)_d\), of dimension \(Q(d)\).
Consequently its cokernel has the fixed fibre dimension
\[
             |\mathcal H|-\bigl(\dim_K A_d(K)-Q(d)\bigr).
\]
Apply Lemma~\ref{lem:ftb-very-dense-projective} to this
finite cokernel.
It is projective.
The exact sequence with target \(R^{\mathcal H}\) splits,
making \(\operatorname{im}E_d\) projective.
The surjection \(A_d(R)\to\operatorname{im}E_d\) then splits.
Equivalently there is a generalized inverse \(s\) with
\[
                                  E_dsE_d=E_d.
 \tag{M10}\label{eq:ftb-evaluation-generalized-inverse}
\]
The projection \(1-sE_d\) has image \(\ker E_d\).
This kernel is finite projective and commutes with
\emph{every} base change, since the displayed splitting
survives tensor product.
The ideal it generates is therefore finitely generated,
and for every field map its extension is exactly the
ideal generated by the new degree-\(d\) evaluation kernel.
This proves the second assertion of
\eqref{eq:ftb-bounded-form-projectivity}; it is not an
unjustified assertion that kernels commute with field
specialization.

At a good prime and any \(e\ge d\), the fibre of
\(A_e(R)/J_e\) now has dimension
\(\dim_KA_e(K)-Q(e)\), by
\eqref{eq:ftb-bounded-cutby}.
Homogeneous-part formation commutes with the field map:
write a polynomial as the sum of its homogeneous parts,
and project to degree \(e\).
The relevant module is a quotient of the finite module
\(A_e(R)\), so it is finite.
Lemma~\ref{lem:ftb-very-dense-projective} makes it
projective.
This proves the first assertion.

Construct the homogeneous projective coordinate scheme
of \(J\), as in the preceding subsection.
Finite generation of \(J\) gives finite presentation
on every standard coordinate chart.
We give the exact flat-chart argument used in the source.
On the \(i\)-th chart its coordinate evaluation has
\(X_i=1\).
Given a finite linear relation
\(\sum_a r_az_a=0\) in this chart, represent the \(z_a\)
by homogeneous forms and homogenize them all to one
degree by powers of \(X_i\).
The homogeneous relation vanishes after localization,
so some additional power \(X_i^k\) puts it in \(J\).
Multiply further by \(X_i^d\).
The relation is now a relation in the projective
module \(A_{d+k+e}(R)/J_{d+k+e}\).
Its flatness expresses it as a trivial relation there:
the representative vector is a linear combination of
vectors of coefficients whose scalar combinations
vanish.
Evaluate these representatives at \(X_i=1\).
The same trivial-relation presentation is obtained for
the original chart relation.
This is the finite-relations criterion for flatness.
Thus every chart, and therefore \(Z/R\), is flat.

Finally fix one \(I_x\) as in the assertion.
The two homogeneous ideals \(I_x\) and \(J_d(xp)\)
have the same parts in all degrees \(e\ge d\).
They consequently have the same projective coordinate
functor on all \(K\)-algebras.
Here is the lower-degree check.
If a homogeneous form \(F\) has degree below \(d\),
multiply it successively by coordinate variables.
Once its degree reaches \(d\), the required vanishing
is known.
If every \(v_iF(v)\) vanishes and the coordinates
\(v_i\) generate the unit ideal, then \(F(v)=0\).
Descending the degree proves the same assertion for
\(F\) itself.
Taking homogeneous parts handles arbitrary forms.
Hence an arbitrary unimodular coordinate vector
vanishes on one ideal if and only if it vanishes on
the other.
The coordinate representability identities give an
isomorphism of the two projective schemes.
On local coordinate frames it preserves the ratios of
the sections, so it gives the actual line-bundle
isomorphism and equality of section coordinates.
Together with the already proved base change of \(J\),
this identifies it with the fibre of \(Z\).
\end{proof}

This proposition is the source's bounded-degree reconstruction,
rather than a Hilbert--Samuel or generic extension substitute.
Its next actual application still requires the source to supply
one degree \(d\), one polynomial \(Q\), and the precise marked
torsion coordinate vectors cutting every permitted theta
fibre.  The analogous graph application needs the corresponding
uniform graph-cut statements.  Those suppliers have not been
asserted merely by proving the algebra above.

\paragraph{Frozen source concordance.}
The very-dense condition, \(\mathrm{dimDeg}\),
\(\mathrm{vanIdeal}\) and \(\mathrm{CutBy}\) are defined in
\srcid{src/Development/AbelianVarietyTheory/Definitions/ModuliSpaces/%
FiniteHomogeneousInf.lean}
(\srcid{1a1beedf08bb1a0f2369c22fa441010a9df8d31b}).
The reduced constant-rank splitting is
\srcid{src/Development/AbelianVarietyTheory/Proofs/FoundationProofs/%
ProjectiveCokerOfIsLocallyConstantFinrankRange.lean}
(\srcid{044befc172ee0d741f991579401dacb0c1aa5164}).
The actual evaluation map, generalized inverse, base-change
kernel and high-degree projectivity are in
\srcid{src/Development/AbelianVarietyTheory/Proofs/ModuliProofs/%
HomogeneousComponentMemVanIdeal.lean}
(\srcid{8098091c209f1f8e587027647f17f7df83182d0c}).
The flat standard-chart calculation is
\srcid{ExistsProjCoordRepresentationHomogeneousIdealFla.lean}
(\srcid{5b0ceb928553b508f4c65e5c8a63dab7f5ceefdc}),
in the same proof directory.
The exact all-fibre assembly and lower-degree argument are
\srcid{ExistsFlatProjectiveThroughPointsGoodFibres.lean}
(\srcid{d1a5d3c35301a311389973b14729a24844b62260}).
These are actual producer bodies, with reducedness,
very density and the same degree/polynomial retained.
